\documentclass[12pt]{article}
\usepackage[T1]{fontenc}
\usepackage[utf8]{inputenc}
\usepackage{lmodern}
\usepackage{textcomp}
\usepackage{enumitem}
\usepackage{geometry}
\geometry{margin=1in}
\begin{document}
\begin{center}
Answer Key - Political Science - Graduate Level Assessment 13 \
\small (Answers are ordered exactly as in the question paper.)
\end{center}

\section*{Multiple Choice}
\begin{enumerate}[label=\arabic*.]
\item \textbf{Answer:} A
\\ \textit{Options:} A) There is no specific referent point.; B) The state and state apparatus.; C) The environment.; D) The people and their collectives.
\\ \textit{Note:} In Critical Security Studies, the concept challenges traditional notions by arguing that security should not be confined to a singular referent object, such as the state or military, but instead emphasizes a broader understanding of security that considers various subjects and contexts.
\item \textbf{Answer:} A
\\ \textit{Options:} A) US banks funded the allies much more than Germany and Austria-Hungary; B) US banks funded Germany and Austria-Hungary much more than the allies; C) Wilson secretly favoured the German position; D) None of the above
\\ \textit{Note:} The correct answer highlights that the substantial financial support provided by US banks to the Allies created a vested economic interest, which ultimately influenced America's neutrality and leaned it towards supporting the Allies in World War I.
\item \textbf{Answer:} D
\\ \textit{Options:} A) Vietnam War.; B) Second World War.; C) First World War.; D) 1991 Gulf War.
\\ \textit{Note:} The 1991 Gulf War exemplifies overwhelming technological superiority, as the coalition forces employed advanced weaponry and precision airstrikes that swiftly incapacitated the Iraqi military, leading to a rapid and decisive victory.
\item \textbf{Answer:} A
\\ \textit{Options:} A) It criticized international organizations, rather than trying to strengthen them; B) It expanded NATO to include former Soviet states; C) It focused on a more personal style of leadership; D) It increased international support for the United States
\\ \textit{Note:} The George W. Bush administration shifted U.S. foreign policy by adopting a unilateral approach that often criticized and sidelined international organizations, such as the United Nations, in favor of pursuing American interests independently, as exemplified by the decision to invade Iraq without broad international support.
\item \textbf{Answer:} B
\\ \textit{Options:} A) To increase the wealth of King George III; B) The growing costs of war with France; C) Anger at America's growing prosperity; D) Pressure from rich merchants
\\ \textit{Note:} The correct answer is based on the fact that the financial burden of the Seven Years' War, particularly its expenses in North America, prompted Britain to seek new revenue sources through taxation of its American colonies to help cover the costs of the conflict and subsequent military presence.
\end{enumerate}

\section*{True / False}
\begin{enumerate}[label=\arabic*.]
\setcounter{enumi}{5}
\item \textbf{Answer:} False
\\ \textit{Options:} A) True; B) False
\\ \textit{Note:} American exceptionalism emphasizes the unique role and values of the United States in the world, often advocating for national sovereignty and unilateral action rather than endorsing a world government as a solution to global governance challenges.
\item \textbf{Answer:} True
\\ \textit{Options:} A) True; B) False
\\ \textit{Note:} The National Security Council is indeed responsible for coordinating American foreign and military policy as it serves as the principal forum for the President to discuss and evaluate national security and foreign policy issues with senior national security advisors and cabinet officials.
\item \textbf{Answer:} False
\\ \textit{Options:} A) True; B) False
\\ \textit{Note:} NATO is primarily a military alliance focused on collective defense and security rather than on economic development, which is more the domain of organizations like the International Monetary Fund (IMF) or the World Bank.
\item \textbf{Answer:} True
\\ \textit{Options:} A) True; B) False
\\ \textit{Note:} The statement is true because the United Nations Security Council is structured with five permanent members-namely the United States, Russia, China, France, and the United Kingdom-who possess veto power over resolutions, while the remaining ten non-permanent members are elected for two-year terms and do not have such veto authority.
\item \textbf{Answer:} True
\\ \textit{Options:} A) True; B) False
\\ \textit{Note:} NSC-68 marked a significant shift in U.S. foreign policy by advocating for a robust military response to the perceived threat of communism on a global scale, thereby militarizing the strategy of containment and promoting the development of the hydrogen bomb as a means to ensure national security and deter Soviet aggression.
\end{enumerate}

\section*{Long Form Response}
\begin{enumerate}[label=\arabic*.]
\setcounter{enumi}{10}
\item \textbf{Answer:} States may be compelled to join alliances like NATO due to a combination of political, economic, and security motivations. Politically, membership can enhance a state's international legitimacy and provide a platform for influence within a collective framework. Economically, alliances can lead to increased trade opportunities and access to resources, as well as financial support for defense initiatives. From a security perspective, joining an alliance like NATO offers collective defense guarantees that deter potential aggressors and foster stability in a volatile geopolitical landscape. Thus, the interplay of these factors-political legitimacy, economic benefits, and enhanced security- illustrates the multifaceted motivations behind a state's decision to pursue alliance membership.
\\ \textit{Note:} correct because it effectively identifies and elaborates on the intertwined political, economic, and security motivations that drive states to join alliances like NATO, emphasizing how such memberships enhance legitimacy, foster trade, and provide collective defense assurances.
\item \textbf{Answer:} The 1994 Human Development Report, published by the United Nations Development Programme (UNDP), was pivotal in redefining the concept of development beyond mere economic growth to include human well-being and capabilities. This shift emphasized that development should focus on enhancing individuals' freedoms and opportunities, thereby influencing contemporary political science discourse on the intersection of human rights, social justice, and policy-making. By introducing the Human Development Index (HDI), the report provided a comprehensive framework for assessing progress, which has since informed national and international policies aimed at improving quality of life. Consequently, policymakers increasingly prioritize human-centric approaches, recognizing that sustainable development hinges on addressing inequalities and empowering individuals. This foundational document has thus shaped both scholarly debates and practical strategies in the realm of human development.
\\ \textit{Note:} correct because it captures the transformative impact of the 1994 Human Development Report in broadening the definition of development to prioritize human well-being and capabilities, which has significantly influenced contemporary political science discussions on human rights, social justice, and policy-making.
\end{enumerate}

\vfill
\noindent\textit{End of Answer Key}
\end{document}